%! TEX root = main.tex

\section{Wave Reflection at an Impedance Jump}
\label{sc:13}

The propagation of a time-harmonic field in a heterogeneous one-dimensional medium is considered, with a discontinuity in the material properties located at mid-domain. The main effect of such discontinuity is the generation of a partially reflected wave due to an impedance mismatch.

\subsection{Problem Setup}

The spatial domain is defined as \(\Omega = (0,1)\) with constant density \(\rho(x) = 1\). The stiffness is assumed piecewise constant and discontinuous at \(x = 0.5\), according to

\begin{equation}
    \mu(x) =
    \begin{cases}
        4, & x \leq 0.5 \\
        1, & x > 0.5
    \end{cases}
    \label{eq:mu_jump}
\end{equation}

No volumetric source term is applied, hence \(f(x) = 0\). The excitation is imposed at the left boundary through a Neumann condition \(\mu(0)u'(0) = g_N\) with \(g_N = 1\), whereas at the right boundary a perfectly absorbing condition is prescribed by setting \(g_A = 0\). This choice prevents spurious reflections at \(x = 1\), so that the observed reflected component is generated only at the internal interface \(x = 0.5\).

\subsection{Mechanical Impedance of the Two Media}

For one-dimensional wave propagation, reflection and transmission at an interface are governed by the characteristic (mechanical) impedance of the medium, defined as

\begin{equation}
    Z = \sqrt{\rho\mu}
    \label{eq:impedance}
\end{equation}

Since \(\rho = 1\) throughout \(\Omega\), the impedance depends only on \(\mu\). In the left subdomain \(x \leq 0.5\) it follows that \(Z_1 = \sqrt{1 \cdot 4} = 2\), while in the right subdomain \(x > 0.5\) one obtains \(Z_2 = \sqrt{1 \cdot 1} = 1\). The corresponding phase velocities are \(c_1 = \sqrt{\mu_1/\rho} = 2\) and \(c_2 = \sqrt{\mu_2/\rho} = 1\) respectively.

\subsection{Reflection Coefficient at the Interface}

When a right-going incident wave traveling in medium 1 reaches the interface with medium 2, continuity of displacement \(u\) and traction \(\mu u'\) implies the formation of both a reflected and a transmitted component. The reflection coefficient for incidence from medium 1 to medium 2 is given by

\begin{equation}
    R = \frac{Z_1 - Z_2}{Z_1 + Z_2}
    \label{eq:reflection_coeff}
\end{equation}

Substituting \(Z_1 = 2\) and \(Z_2 = 1\) yields
\[
R = \frac{2 - 1}{2 + 1} = \frac{1}{3}
\]
Therefore, a partial reflection is expected at \(x = 0.5\) with positive sign, i.e., without phase inversion.

\subsection{Expected Qualitative Behavior in Amplitude and Phase}

Distinct qualitative behaviors are expected on the two sides of the interface. In the left region \(x < 0.5\), the superposition of the incident and reflected waves produces an interference pattern, corresponding to a partially standing-wave regime. As a consequence, oscillations are expected in both the amplitude profile \(|u(x)|\) and the phase \(\arg(u(x))\), which will be irregular and non-monotone. In the right region \(x > 0.5\), the transmitted wave propagates toward \(x = 1\). Due to the perfectly absorbing condition \(g_A = 0\), no reflection is generated at the right boundary and the solution resembles a single traveling wave. Hence, the amplitude is expected to be comparatively smooth and the phase to grow linearly and monotonically with slope \(\omega/c_2 = \omega\). At the interface \(x = 0.5\), a discontinuity in the spatial derivative of the phase is expected, reflecting the change in phase velocity from \(c_1 = 2\) to \(c_2 = 1\): since the phase slope of a traveling wave is \(\omega/c\), the slope in the right subdomain is twice that in the left subdomain for a purely traveling wave.

\subsection{Numerical Results}

The problem is solved numerically using the MATLAB solver described in chapter \ref{sc:11}, with \(N = 200\) elements to ensure sufficient resolution across both subdomains. The stiffness coefficient is implemented as a piecewise function evaluated at the quadrature nodes of each element. Figure \ref{fig:ampl_phase} reports the amplitude \(|u(x)|\) and the unwrapped phase \(\arg(u(x))\) of the computed solution.

\begin{figure}[!htb]
    \centering
    \subfloat[Amplitude \(|u(x)|\) of the numerical solution. The oscillatory pattern in the left subdomain (\(x < 0.5\)) is the signature of the standing-wave component generated by the partial reflection at the interface, while the smoother decay in the right subdomain (\(x > 0.5\)) corresponds to the purely transmitted wave]{\includegraphics[width=0.90\textwidth]{00_Images/HW1/3.1_amplitude.png}}
    \\
    \subfloat[Unwrapped phase \(\arg(u(x))\) of the numerical solution. The irregular, non-monotone behavior for \(x < 0.5\) is due to the interference between the incident and reflected waves, while the linear monotone growth for \(x > 0.5\) is consistent with a single right-going traveling wave. The kink at \(x = 0.5\) reflects the change in phase velocity between the two media]{\includegraphics[width=0.90\textwidth]{00_Images/HW1/3.2_phase.png}}
    \caption{Amplitude and unwrapped phase of the FE solution for the impedance-jump problem with \(\mu\) as in equation \eqref{eq:mu_jump}, \(\rho = 1\), \(g_N = 1\), \(g_A = 0\) and \(\omega = 5\pi\)}
    \label{fig:ampl_phase}
\end{figure}

\subsection{Physical Interpretation}

The numerical results are in full agreement with the analytical predictions derived in the previous sections. Three distinct regions can be identified in the solution.

\paragraph{Left subdomain \(x < 0.5\) : partial standing wave}
The amplitude profile \(|u(x)|\) exhibits oscillations with alternating local maxima and minima, which are the signature of the interference between the rightward-traveling incident wave and the leftward-traveling reflected wave generated at the interface. The phase is irregular and non-monotone, consistently with the superposition of two counter-propagating components. The reflection coefficient \(R = 1/3\) implies that approximately \(11\%\) of the incident energy is reflected back into the left subdomain, producing a moderately pronounced standing-wave pattern.

\paragraph{Right subdomain \(x > 0.5\) : traveling wave}
Beyond the interface, the amplitude is smooth and the phase grows linearly with slope \(\omega/c_2 = \omega\), confirming the presence of a single right-going traveling wave. The absorbing boundary condition at \(x = 1\) effectively prevents any reflection from the right end, so no additional interference pattern appears in this region.

\paragraph{Interface \(x = 0.5\) : phase velocity change}
A clear kink in the phase is visible at the interface. Since the phase slope of a traveling wave equals \(\omega/c\), the ratio of the slopes on the two sides of the interface is
\[
\frac{(\omega/c_2)}{(\omega/c_1)} = \frac{c_1}{c_2} = \frac{2}{1} = 2
\]
confirming that the phase advances twice as fast per unit length in the right medium, consistently with \(c_2 = 1 < c_1 = 2\).